\chapter{Удиртгал}

\section{Системийн онол арга зүйн судалгаа}

\subsection{Системийн онолын үндэс}
Өнөөгийн зах зээл дээрх мессежийн аппликейшнүүд нь хэрэглэгчийн туршлагыг (UX) сайжруулахын тулд төрөл бүрийн шүүлтүүр болон хэрэглэгчийн харилцах цонхны шийдлүүдийг ашиглаж байна. Судалгааны хүрээнд дараах 4 гол системийг сонгон авч шинжлэв.

\subsection{Ижил төстэй системүүдийн судалгаа}
\begin{itemize}
    \item \textbf{Google Messages}: Android-ийн үндсэн систем бөгөөд спам шүүлтүүрийн ``Machine Learning'' алгоритм сайтай боловч хэрэглэгч өөрөө шүүлтүүрийн дүрмийг (Regex) гараар тохируулах боломжгүй.
    \item \textbf{Samsung Messages}: Samsung төхөөрөмжүүдийн нугалдаг төхөөрөмжүүдийн хэрэглэгчийн харилцах цонхыг (Two-pane layout) хамгийн сайн хэрэгжүүлсэн хувилбар. Гэвч зөвхөн Samsung-ийн экосистемд хаалттай ажилладаг.
    \item \textbf{Pulse SMS}: Олон төхөөрөмж дээр синхрончлох болон хавтаслах системтэй боловч автомат ангилал нь зөвхөн энгийн түлхүүр үгээр хязгаарлагддаг.
    \item \textbf{Textra SMS}: Хэрэглэгчийн харилцах цонхны өнгө, дизайны өргөн сонголттой боловч орчин үеийн нугалдаг дэлгэцийн (Adaptive UI) оновчлол сул байна.
\end{itemize}

\subsection{Системүүдийн харьцуулсан шинжилгээ}
Дээрх системүүдийн функционал болон хэрэглэгчийн харилцах цонхны боломжуудыг өөрийн хөгжүүлж буй системтэй Хүснэгт \ref{table:comparison}-д харьцуулан харуулав.

\begin{table}[H]
\centering
\caption{Ижил төстэй системүүдийн харьцуулалт}
\label{table:comparison}
\begin{tabular}{|l|c|c|c|c|c|}
\hline
\textbf{Үзүүлэлтүүд} & \textbf{Google} & \textbf{Samsung} & \textbf{Pulse} & \textbf{Textra} & \textbf{Fossify (Enhanced)} \\ \hline
Regex Matching & \xmark & \xmark & \xmark & \xmark & \cmark \\ \hline
Two-pane Layout & \cmark & \cmark & \cmark & \xmark & \cmark \\ \hline
Open Source & \xmark & \xmark & \xmark & \xmark & \cmark \\ \hline
Custom Folders & \xmark & \xmark & \cmark & \xmark & \cmark \\ \hline
Swipe Actions & \cmark & \cmark & \cmark & \cmark & \cmark \\ \hline
\end{tabular}
\end{table}

\subsection{Судалгааны үр дүн ба зөрүүний шинжилгээ}
Харьцуулсан шинжилгээнээс үзэхэд дараах дүгнэлтүүд гарч байна:
\begin{enumerate}
    \item \textbf{Шүүлтүүрийн дутагдал}: Одоогийн системүүд спамыг автоматаар хаадаг боловч хэрэглэгч өөрийн хүссэнээр (жишээ нь: банкны гүйлгээг ``Finance'' хавтсанд) Regex ашиглан нарийн ангилах боломж хязгаарлагдмал байна.
    \item \textbf{Хэрэглэгчийн харилцах цонхны үл нийцэл}: Google болон Samsung-аас бусад аппликейшнүүд нугалдаг дэлгэцийн ``State Continuity'' болон ``Large Screen'' оновчлолыг орхигдуулсан байна.
    \item \textbf{Нээлттэй эхийн хэрэгцээ}: Хэрэглэгчийн хувийн зурваст ханддаг систем нь ``Open Source'' байх нь аюулгүй байдал болон хяналтын хувьд чухал ач холбогдолтой юм.
\end{enumerate}

Энэхүү төслөөр Fossify Messages-ийн нээлттэй эх дээр тулгуурлан тогтмол илэрхийлэл ашигласан шүүлтүүр болон нугалдаг дэлгэцийн оновчлолыг нэгтгэж, зах зээлд байгаа эдгээр дутагдалтай талуудыг (Gap) шийдвэрлэхийг зорьж байна.

\section{Төслийн үндэслэл}
Сүүлийн жилүүдэд Монгол улсын хэмжээнд цахим шилжилт эрчимтэй явагдаж, төрийн болон хувийн хэвшлийн үйлчилгээний мэдэгдэл, нэг удаагийн код (OTP) болон банкны гүйлгээний мэдээлэл мессежээр дамжих нь эрс нэмэгдсэн. Практик ажиглалтаас үзэхэд ихэнх хэрэглэгчийн
хүлээн авч буй зурвасын дийлэнх хэсгийг автомат системээс илгээсэн үйлчилгээний мэдэгдлүүд эзэлж байна. Эдгээр зурвасууд нь хувийн харилцааны мессежүүдтэй нэг жагсаалтад холилдон харагдах нь хэрэглэгчийн хэрэгцээт мэдээллээ хайх хугацааг ихэсгэж, мэдээллийн ачаалал үүсгэх гол хүчин зүйл болж байна.

\section{Асуудлын тодорхойлолт}
Одоогийн мессежийн аппликейшнуудын хувьд зурвасуудыг агуулгаар нь таньж, автоматаар хавтаслах механизм дутагдалтай байгаагаас дараах асуудлууд үүсэж байна:
\begin{itemize}
    \item \textbf{Эмх цэгцгүй байдал}: Санхүү, сурталчилгаа, хувийн харилцааны зурвасууд нэг дор холилдож харагддаг.
    \item \textbf{Автоматжуулалт сул}: Давтамжтай ирдэг мэдэгдлүүдийг (жишээ нь: 150000 дугаараас ирэх гүйлгээний мэдээлэл) таних уян хатан шүүлтүүр байхгүй.
    \item \textbf{Төхөөрөмжийн үл нийцэл}: Дэлгэддэг болон өргөн дэлгэцтэй төхөөрөмжүүд дээрх том талбайг үр ашигтай ашиглах хэрэглэгчийн харилцах цонхны шийдэл хангалтгүй (Fossify Messages- системийн хувьд).
\end{itemize}

\textbf{Оролцогч талууд:}
\begin{enumerate}
    \item \textbf{Үндсэн хэрэглэгчид}: Өдөр тутам их хэмжээний үйлчилгээний мэдэгдэл хүлээн авдаг, мэдээллээ цэгцлэх сонирхолтой иргэд.
    \item \textbf{Захиалагч/Хөгжүүлэгч}: Нээлттэй эх (Fossify) дээр суурилсан шийдлийг оновчтой болгож буй баг.
    \item \textbf{Төхөөрөмж үйлдвэрлэгчид}: Нугалдаг болон таблет төхөөрөмжийн хэрэглээний тасралтгүй байдлыг дэмжих хэрэгцээтэй тал.
\end{enumerate}

\section{Төслийн зорилго}
Төслийн гол зорилго нь нээлттэй эхийн Fossify Messages \cite{fossify_messages} '' аппликейшнд суурилан хэрэглэгчийн ирж буй зурвасыг автоматжуулсан дүрмийн дагуу хавтаслах систем болон орчин үеийн уян хатан хэрэглэгчийн харилцах цонхны шийдлийг хөгжүүлэхэд оршино. Төслийн зорилгыг SMART зарчмын~ \cite{doran_smart} дагуу дараах байдлаар тодорхойлж байна:

\begin{itemize}
    \item \textbf{Specific (Тодорхой байдал)}: Хэрэглэгчийн тохируулсан түлхүүр үг болон тогтмол илэрхийлэл дээр суурилсан мессеж ангилах алгоритмыг хөгжүүлж, ирсэн зурвасыг тохирох хавтсанд хүний оролцоогүйгээр автоматаар ангилан хуваарилдаг механизмыг хэрэгжүүлэх.
    \item \textbf{Measurable (Хэмжигдэхүйц байдал)}: Автомат ангиллын алгоритмыг үнэлэх тестүүдийг бэлдэж, урьдчилан тодорхойлсон дүрмүүдэд нийцэж буй ирсэн мессежийн 100\%-ийг алдаагүй, зөв хавтсанд оруулах нөхцөлийг хангах.
    \item \textbf{Achievable (Хэрэгжихүйц байдал)}: Хэрэглэгчийн хэрэглэгчийн харилцах цонхны үйлдлүүдийг оновчилж, шудрах болон удаан дарах хөдөлгөөнүүдийг нэгтгэн, гол үйлдлүүдийг хамгийн ихдээ 2 алхам дотор гүйцэтгэх урсгалыг бүрдүүлэх; мөн нугалдаг төхөөрөмжид зориулсан хоёр баганат харагдацын тасралтгүй ажиллагааг хангах.
    \item \textbf{Relevant (Ач холбогдолтой байдал)}: Энэхүү шийдэл нь Монгол улсын цахим шилжилтийн хүрээнд хэрэглэгчдэд практик үнэ цэнэтэй бөгөөд нээлттэй эхийн хөгжүүлэлтийн нийгэмлэгт хувь нэмэр болно.
    \item \textbf{Time-bound (Хугацааны хязгаарлалт)}: Хөгжүүлэлтийн бүх шатыг дипломын ажлын календарчилсан төлөвлөгөөний дагуу гүйцэтгэж, төгсгөлийн хамгаалалтаас өмнө бүрэн ажиллагаатай хувилбарыг бэлэн болгох.
\end{itemize}

\section{Хамрах хүрээ ба хязгаарлалт (Scope and Limitations)}
Төслийн хүрээнд Android систем \cite{android_dev_docs} дээрх зурвас ангилах тогтмол илэрхийллийн хөдөлгүүр болон дасан зохицох хэрэглэгчид харагдах хэсгийн модулийг хөгжүүлнэ.
\begin{itemize}
    \item \textbf{Хамрах хүрээ}: Богино мессеж, локал өгөгдлийн сангийн шүүлтүүр,түүнтэй холбогдох хэрэглэгчийн харилцах хэсэг болон нугалдаг төхөөрөмжийн хоёр баганат бүтэц, шилжилтын тасралтгүй байдал.
    \item \textbf{Хязгаарлалт}: Систем нь зөвхөн төхөөрөмж дээр локал байдлаар ажиллах бөгөөд Cloud синхрончлол болон RCS (Rich Communication Services) функцуудыг хамрахгүй.
\end{itemize}
